% =========================================================== Q4
\section*{Q4: R-GCN Implementation (6 pts)}

\definecolor{pTwoKw}{RGB}{0,92,197}
\definecolor{pTwoCmt}{RGB}{106,115,125}
\definecolor{pTwoStr}{RGB}{3,47,98}

\lstdefinestyle{pTwoPyStyle}{%
  language=Python,
  basicstyle=\ttfamily\scriptsize,
  keywordstyle=\color{pTwoKw}\bfseries,
  commentstyle=\color{pTwoCmt}\itshape,
  stringstyle=\color{pTwoStr},
  showstringspaces=false,
  columns=fullflexible,
  keepspaces=true,
  breaklines=true,
  breakautoindent=true,
  xleftmargin=1.0em,
  numbers=none,
  tabsize=4,
}

\begin{lstlisting}[style=pTwoPyStyle]
class MyRGCNConv(MessagePassing):
    """R-GCN: one weight matrix per relation, plus a self-connection."""

    def __init__(self, in_channels, out_channels, num_relations):
        # aggr="mean" IS the 1/c_{i,r} of the formula, taking
        # c_{i,r} = |N_r(i)|: the example Schlichtkrull et al.'s Eq. (2)
        # itself offers. PyG's mean is 0, not NaN, when |N_r(i)| = 0, so
        # no empty neighbourhood needs a special case.
        super().__init__(aggr="mean")
        self.in_channels = in_channels
        self.out_channels = out_channels
        self.num_relations = num_relations
        self.root = nn.Linear(in_channels, out_channels, bias=True)  # W_0
        # W_r: one full [in, out] matrix per relation. Basic R-GCN, so
        # neither basis nor block-diagonal decomposition.
        self.rel = nn.ModuleList([
            nn.Linear(in_channels, out_channels, bias=False)
            for _ in range(num_relations)])
        self.reset_parameters()

    def reset_parameters(self):
        # Safe to call twice, and the harness does call it twice.
        super().reset_parameters()
        self.root.reset_parameters()
        for lin in self.rel:
            lin.reset_parameters()

    def forward(self, x, edge_index, edge_type):
        out = self.root(x)                             # W_0 h_i
        n = x.size(self.node_dim)
        for r in range(self.num_relations):
            ei_r = edge_index[:, edge_type == r]       # relation r only
            if ei_r.size(1) == 0:
                continue
            # c_{i,r} does not depend on j and W_r is linear, so
            # sum_j (1/c_{i,r}) W_r h_j == W_r (mean_j h_j): transform
            # first, propagate second. size=(n, n) keeps the result at the
            # full node count where relation r touches only some nodes.
            out = out + self.propagate(ei_r, x=self.rel[r](x), size=(n, n))
        return out

    def message(self, x_j):
        return x_j
\end{lstlisting}

\noindent
\texttt{self.rel[r]} is $W_r$: one full $\mathrm{in}\times\mathrm{out}$ matrix per
relation, with neither of the decompositions Schlichtkrull et al.~\cite{p2:rgcn} add as
regularisers~--- the \emph{basic} version. \texttt{self.root} is $W_0$, applied
to $h_i$ outside any propagation.

I took $c_{i,r}=|N_r(i)|$, the example their Eq.~(2) offers, which is exactly
\texttt{aggr="mean"}. Guarding it is not optional: all 15{,}466 nodes have
$|N_r(i)|=0$ for some $r$; PyG's mean divides by \texttt{count.clamp(min=1)},
returning $0$ rather than NaN.

\texttt{edge\_type} never reaches \texttt{propagate}: \texttt{forward} slices
\texttt{edge\_index} by relation and propagates once per subset, so relation $r$'s
edges only meet $W_r$ and the stub's
\texttt{message(self,\,x\_j)} is untouched. $W_r$ is applied before propagating rather
than after~--- $c_{i,r}$ does not depend on $j$ and $W_r$ is linear, so
$\sum_{j} c_{i,r}^{-1} W_r h_j = W_r\!\left(\operatorname{mean}_{j} h_j\right)$~---
transforming 15{,}466 rows instead of 3{,}014{,}361 messages. \texttt{size=(n,\,n)}
keeps every relation's output at the full node count.

\medskip
\noindent
Cell 28's \texttt{train\_and\_eval\_hetero("R-GCN")}, run unmodified, selected
epoch~13:

\begin{center}
\begin{tabular}{lcc}
\toprule
 & MRR & Hit@1 \\
\midrule
validation & 0.2291 & 0.0782 \\
test       & 0.2229 & 0.0714 \\
\bottomrule
\end{tabular}
\end{center}

\noindent
Random tie-breaking over 21 candidates already scores MRR $0.174$, so the model barely
beats chance. Cell 26's features come from gensim's \texttt{Word2Vec} at
\texttt{workers=4}, not reproducible across processes; over five runs validation MRR was
$0.190\pm0.020$ and test MRR $0.182\pm0.023$. The layer is not the cause:
\texttt{torch\_geometric.nn.RGCNConv} in the same harness scores $0.208$, and my
\texttt{forward} matches it to $4\times10^{-7}$ on all 6{,}067{,}650 edges.
